\cotefile{64}
\scanpage{116}
% archive foliation: 115
soit
\[
\tag{61}
\sigma_i z=\frac{\zeta_i z+\overbrace{\zeta_\lambda\zeta_\mu-\zeta_i(\zeta_\lambda+\zeta_\mu)}^{b_i}}{z-\zeta_i}.
\]
On peut transformer $b_i$ en utilisant les fonctions symétriques élémentaires $(-1)^\alpha a_\alpha$ des $\zeta$,
\[
b_i=2\underbrace{\zeta_\lambda\zeta_\mu}_{-a_3/\zeta_i}-\underbrace{(\zeta_i\zeta_\lambda+\zeta_\lambda\zeta_\mu+\zeta_\mu\zeta_i)}_{a_2}=-(a_2+2a_3/\zeta_i)
\]
donc on trouve
\[
\tag{62}
\sigma_i z=\frac{\zeta_i z-(a_2+2a_3/\zeta_i)}{z-\zeta_i}=\frac{\zeta_i^2z-(a_2\zeta_i+2a_3)}{\zeta_i z-\zeta_i^2}
\]
\margin{expressions valables du moins si $\zeta_i\ne0$, sinon on convient bien sûr que $a_3/\zeta_i$ \unsure{signifie} $\unsure{-}\zeta_\lambda\zeta_\mu$. Mais quitte à modifier \unsure{les $a_i$}, on \unsure{peut même supposer} les $\zeta_i\ne0$\ldots}
Les \inserted{div. $\Delta_i$ des} points fixes de $\sigma_i$ est l'ens. div. des solutions de l'équation en $z_i$
\[
\tag{63}
z_i^2-2\zeta_i z_i+(a_2+2a_3/\zeta_i)=0
\]
dont les solutions sont données par
\[
\tag{64}
z_i=\zeta_i\pm\sqrt{\underbrace{\zeta_i^2-a_2-2a_3/\zeta_i}_{\deleted{\unsure{\zeta}}}}
\]
L'expression sous le signe $\sqrt{\phantom{x}}$ s'écrit aussi
\[
\frac{1}{\zeta_i}(\zeta_i^3-a_2\zeta_i-2a_3)
\]
ou encore, compte tenu de (60)
\[
\zeta_i^3=-a_1\zeta_i^2-a_2\zeta_i-a_3
\]
\deleted{ce radical} cette expression devient
\[
-\frac{1}{\zeta_i}(a_1\zeta_i^2+2a_2\zeta_i+3a_3)
\]
donc (64) s'écrit aussi
\[
\tag{\textcircled{65}}
\left\{\begin{aligned}
(z_i-\zeta_i)^2&=-\frac{1}{\zeta_i}(a_1\zeta_i^2+2a_2\zeta_i+3a_3)\\
\text{i.e.}\quad z_i&=\zeta_i\pm\sqrt{-\frac{1}{\zeta_i}(a_1\zeta_i^2+2a_2\zeta_i+3a_3)}
\end{aligned}\right.
\]
\margin{Si $\zeta_i=0$, on \deleted{\unsure{peut écrire ainsi}} $-(a_1\zeta_i+2a_2+\unsure{3\zeta_\lambda\zeta_\mu|\zeta_\lambda\zeta_\mu})$ l'expr. sous le radical\ldots}
Les douze points d'ordre $4$ de $E$ sont donc les points au dessus des six points $z_i$, donc \deleted{définis par} \inserted{en tout}
